Fig.~\ref{fig:resolution}a shows how detection falls with grid spacing. Interpolation scores {{VAL:sw_2_BICUBIC}} at 2\,m, {{VAL:sw_10_BICUBIC}} at 10\,m, {{VAL:sw_30_BICUBIC}} at 30\,m and {{VAL:sw_60_BICUBIC}} at 60\,m. Fig.~\ref{fig:resolution}b relates the fall to face size. At every spacing, boosting finds nearly all faces longer than four grid cells ({{VAL:rec30_GBM_4}} of faces 4 to 8 cells long at 30\,m, {{VAL:rec60_GBM_4}} at 60\,m), but at 30 and 60\,m far fewer faces shorter than half a cell ({{VAL:rec30_GBM_0}} at 30\,m, {{VAL:rec60_GBM_0}} at 60\,m). Longer patches touch more cells and so have more chances to be flagged, which accounts for part of this rise. By count most patches are short: {{VAL:lab30_registered_shortcount}} are shorter than one 30\,m cell, but they hold only {{VAL:lab30_registered_shortarea}} of the face posts, and with a hysteresis rule that joins fragments of one face the shares are {{VAL:lab30_H_shortcount}} and {{VAL:lab30_H_shortarea}} (post hoc).
